\documentclass[12pt,a4paper]{article}
\usepackage[margin=2.5cm]{geometry}
\usepackage{amsmath}
\usepackage{hyperref}
\usepackage{parskip}
\usepackage{booktabs}

\title{\textbf{Data Collection \& Cleaning}\\[0.3em]
\large Part V(a) --- Modelling Report\\
Asymmetric Sentiment Effects in Earnings Calls}
\author{NLP Course --- Group Project}
\date{}

\begin{document}
\maketitle

\section{Overview}

The dataset combines two transcript corpora with daily stock price data from Yahoo Finance.
The goal was to get enough events to run a meaningful event study while keeping
the sample homogeneous enough that a single sentiment-return model makes sense.

\section{Transcript Sources}

\subsection{Thomson Reuters StreetEvents (Primary)}

The primary source is a set of 188 earnings call transcripts for 10 NASDAQ-listed
technology firms covering 2016--2020.
Tickers: AAPL, AMZN, GOOG, META, MSFT, NFLX, NVDA, TSLA, CRM, AMD.
Files follow the naming convention \texttt{YYYY-Mon-DD-TICKER.txt}.

\subsection{Kaggle cleaned\_ECTs (Supplementary)}

To extend coverage and date range, we added transcripts from the publicly available
Kaggle \texttt{cleaned\_ECTs} corpus.
This adds 9 technology-adjacent firms including IBM, Oracle, Salesforce, and Accenture,
extending the date range back to 2007.
Files follow the convention \texttt{YYYY\_Qn\_ticker\_processed.txt}.

\subsection{Combined Dataset}

After filtering both sources to technology and technology-adjacent firms only
(to minimize cross-sector heterogeneity) and removing any duplicates,
the combined corpus contains \textbf{620 events} across \textbf{19 tickers}
spanning 2007--2025.

\begin{center}
\begin{tabular}{ll}
\toprule
Attribute & Value \\
\midrule
Total events & 620 \\
Unique tickers & 19 \\
Date range & 2007--2025 \\
Training window & 2016--2018 \\
Test window & 2019--2020 \\
\bottomrule
\end{tabular}
\end{center}

\section{Market Price Data}

Daily adjusted closing prices are downloaded from Yahoo Finance via \texttt{yfinance}
for each of the 19 tickers plus their benchmark indices:
\texttt{\^{}IXIC} (NASDAQ Composite) for NASDAQ-listed firms and
\texttt{\^{}GSPC} (S\&P 500) for NYSE-listed firms.
We download data from 2006 onward to ensure sufficient pre-event history for the
estimation window.

\section{Event Trading Day Assignment}

A critical cleaning step is assigning each transcript to the correct trading day.
An earnings call held at 10 PM Tuesday affects Wednesday's open; a call at 2 PM
is priced the same day.  We parse each transcript header for a timestamp and apply:

\begin{itemize}
  \item Call before 4:00 PM ET on a trading day $\rightarrow$ event day = same calendar day.
  \item Call at/after 4:00 PM ET, on a weekend, or on a market holiday
        $\rightarrow$ event day = next business day.
\end{itemize}

If no header timestamp is found, we fall back to noon on the calendar date of the
file (same-day assignment).  This rule is implemented in
\texttt{src/nasdaq\_nlp/data/metadata/\_\_init\_\_.py}.

\section{Data Quality}

Events are dropped if the stock or benchmark has fewer than 60 valid trading days
in the estimation window $[-120, -20]$ relative to the event day.
This removes transcripts near IPO dates or during trading halts.
The final 620 events all have complete estimation windows.

\section{Code Reference}

\begin{itemize}
  \item \texttt{src/nasdaq\_nlp/data/loader/} --- transcript scanning for both sources.
  \item \texttt{src/nasdaq\_nlp/data/metadata/} --- header parsing and event day assignment.
  \item \texttt{src/nasdaq\_nlp/data/market/} --- price download and return computation.
  \item Entry point: \texttt{make pipeline} or \texttt{notebooks/01\_data\_pipeline.ipynb}.
\end{itemize}

\end{document}
